\chapter{重定向后的交互强化}

\section{直接使用重定向结果的缺陷}
尽管本方法重定向结果已经达到较高精度和平滑度，但在实际应用中仍存在不足之处：

首先，现实生活中，人在抓取物体的过程中，通常会根据物体来预测大致抓取力，而在本方法的灵巧手抓取物体的过程中，
抓取力在未接触物体时与加速度相关，在与物体接触时，则受到电机最大力矩和物体表面性质影响。
而在此过程中有可能在接触时，就因加速度过大，导致灵巧手的抓取力过大，进而损坏物体。
而一直让灵巧手保持低速运行又显然是一个过于保守的策略。
因此，我们认为在对一个物体进行抓取前，应该先大致预测抓取力，
并在抓取前调整灵巧手电机速度，使得抓取力与预测值一致。
具体的，我们认为可以通过视觉语言模型来预测抓取力，并在即将抓取时，将灵巧手电机速度调整以匹配预测抓取力。

其次，本方法仅根据人手位置生成灵巧手角度，没有考虑灵巧手手指与目标之间的距离等交互信息。
特别是涉及物体表面性质的信息，例如对于易碎或者弹性物体，抓取力可能需要根据物体的材质进行调整，而不仅仅是根据人手的动作来确定。
如果人手在遥操作的过程中，对应灵巧手动作将要穿透物体表面，那么此时就需要根据触觉信息或者关节力矩信息进行修正，
避免抓取动作破坏物体。


\section{基于VLM的目标抓取力预测}

\section{基于力矩反馈的强化学习}

